% GENERATED FILE. Edit canonical lessons, then run npm run generate:books.
% canonical-source-hash: fnv1a64:a5ec5bab73d7369f
% canonical-lessons: LA-C372-auctor, LA-C372-impulsus, LA-C372-efficio, LA-C372-efficax, LA-C372-excito, LA-R372-review-365-372

\chapter{What Sets Things Going}
\label{ch:la-what-sets-things-going}
\hlchaptermodality{372}

\begin{chapteropening}
\textbf{By the end of this chapter:} \emph{I can name the person behind an act and the impulse behind it, and say what brings a result about and what works.}

Complete the last lesson of chapter 372: I can name the person behind an act and the impulse behind it, and say what brings a result about and what works.
\end{chapteropening}

\section[auctor, auctōris]{auctor, auctōris — an originator, an author; a backer}

\label{lesson:LA-C372-auctor}



\begin{quote}
Before the new one: say the Latin for to read aloud, then the Latin for rashly.
\end{quote}

\subsection*{You'll want to know: auctor, auctōris}
\textbf{auctor, auctōris} — \textquotedblleft{}an originator, an author; a backer\textquotedblright{}.

\begin{quote}
Quis huius cōnsiliī auctor fuit?
\end{quote}

\begin{quote}
\emph{Who was the author of this plan?}
\end{quote}

\emph{auctor, auctōris} m. (sometimes f.). The one who makes a thing grow: a founder, a writer, the backer of a law or of a sale. English \emph{author}.

\begin{cousinweb}
\emph{auctor} is built on \emph{augeō} \textquotedblleft{}to increase, to make grow\textquotedblright{} (EtymDB; Lewis \& Short), with the ending \emph{-tor} of the one who does it. English \emph{author} comes from it through French.
\end{cousinweb}

\begin{grammarlens}[title={What you've built}]
One more word for telling, explaining and describing at B1.
\end{grammarlens}

\subsection*{Your turn}
\begin{itemize}
\raggedright
  \item \emph{Say it:} \emph{auctor, auctōris}
  \item \emph{Say it:} \emph{auctor, auctōris}, once more
  \item \emph{Say it:} a sentence of your own with \emph{auctor, auctōris}
  \item \emph{Recall:} say the Latin for to read aloud, then the Latin for rashly, then say \emph{auctor, auctōris} again
\end{itemize}

\begin{tcolorbox}[breakable,colback=teal!4,colframe=teal!35!black,title={Before you move on}]
What does auctor, auctōris mean? (\textquotedblleft{}An originator, an author; a backer\textquotedblright{}.) Say it once more.
\end{tcolorbox}
\section[impulsus, impulsūs]{impulsus, impulsūs — a push; an impulse; instigation}

\label{lesson:LA-C372-impulsus}



\begin{quote}
Before the new one: say the Latin for rashly, then the Latin for an originator.
\end{quote}

\subsection*{You'll want to know: impulsus, impulsūs}
\textbf{impulsus, impulsūs} — \textquotedblleft{}a push; an impulse; instigation\textquotedblright{}.

\begin{quote}
Impulsū amīcōrum hoc fēcit.
\end{quote}

\begin{quote}
\emph{He did this at his friends' urging.}
\end{quote}

\emph{impulsus, impulsūs} m., built on \emph{impellō}, to push on (\emph{in} + \emph{pellō}). \emph{Impulsū alicuius}, at someone's instigation. English \emph{impulse}.

\begin{cousinweb}
\emph{impulsus} is the noun built on \emph{impellō} \textquotedblleft{}to push on\textquotedblright{} (Lewis \& Short), from \emph{in} and \emph{pellō} \textquotedblleft{}to drive\textquotedblright{}. English \emph{impulse} comes from it.
\end{cousinweb}

\begin{grammarlens}[title={What you've built}]
One more word for telling, explaining and describing at B1.
\end{grammarlens}

\subsection*{Your turn}
\begin{itemize}
\raggedright
  \item \emph{Say it:} \emph{impulsus, impulsūs}
  \item \emph{Say it:} \emph{impulsus, impulsūs}, once more
  \item \emph{Say it:} a sentence of your own with \emph{impulsus, impulsūs}
  \item \emph{Recall:} say the Latin for rashly, then the Latin for an originator, then say \emph{impulsus, impulsūs} again
\end{itemize}

\begin{tcolorbox}[breakable,colback=teal!4,colframe=teal!35!black,title={Before you move on}]
What does impulsus, impulsūs mean? (\textquotedblleft{}A push; an impulse; instigation\textquotedblright{}.) Say it once more.
\end{tcolorbox}
\section[efficiō, efficere]{efficiō, efficere — to bring about, to cause; to make}

\label{lesson:LA-C372-efficio}



\begin{quote}
Before the new one: say the Latin for an originator, then the Latin for a push.
\end{quote}

\subsection*{You'll want to know: efficiō, efficere}
\textbf{efficiō, efficere} — \textquotedblleft{}to bring about, to cause; to make\textquotedblright{}.

\begin{quote}
Pluvia effēcit ut domī manērēmus.
\end{quote}

\begin{quote}
\emph{The rain made us stay at home.}
\end{quote}

\emph{efficiō, efficere, effēcī, effectum}. \emph{Efficiō ut} + subjunctive: I bring it about that. Built on \emph{faciō}. English \emph{efficient}.

\begin{cousinweb}
\emph{efficiō} is built from \emph{ex} \textquotedblleft{}out\textquotedblright{} and \emph{faciō} \textquotedblleft{}to make\textquotedblright{} (EtymDB; Lewis \& Short): to make a thing until it is done. English \emph{efficient} comes from its present participle.
\end{cousinweb}

\begin{grammarlens}[title={What you've built}]
One more word for telling, explaining and describing at B1.
\end{grammarlens}

\subsection*{Your turn}
\begin{itemize}
\raggedright
  \item \emph{Say it:} \emph{efficiō, efficere}
  \item \emph{Say it:} \emph{efficiō, efficere}, once more
  \item \emph{Say it:} a sentence of your own with \emph{efficiō, efficere}
  \item \emph{Recall:} say the Latin for an originator, then the Latin for a push, then say \emph{efficiō, efficere} again
\end{itemize}

\begin{tcolorbox}[breakable,colback=teal!4,colframe=teal!35!black,title={Before you move on}]
What does efficiō, efficere mean? (\textquotedblleft{}To bring about, to cause; to make\textquotedblright{}.) Say it once more.
\end{tcolorbox}
\section[efficāx, efficācis]{efficāx, efficācis — effective, efficient}

\label{lesson:LA-C372-efficax}



\begin{quote}
Before the new one: say the Latin for a push, then the Latin for to bring about.
\end{quote}

\subsection*{You'll want to know: efficāx, efficācis}
\textbf{efficāx, efficācis} — \textquotedblleft{}effective, efficient\textquotedblright{}.

\begin{quote}
Hoc remedium efficāx est.
\end{quote}

\begin{quote}
\emph{This remedy is effective.}
\end{quote}

\emph{efficāx, efficācis}: one ending in the nominative for all genders. English \emph{efficacious}.

\begin{cousinweb}
\emph{efficāx} is built on \emph{efficiō} \textquotedblleft{}to bring about\textquotedblright{} (EtymDB; Lewis \& Short), with the ending \emph{-āx} of a lasting tendency. English \emph{efficacious} comes from it.
\end{cousinweb}

\begin{grammarlens}[title={What you've built}]
One more word for telling, explaining and describing at B1.
\end{grammarlens}

\subsection*{Your turn}
\begin{itemize}
\raggedright
  \item \emph{Say it:} \emph{efficāx, efficācis}
  \item \emph{Say it:} \emph{efficāx, efficācis}, once more
  \item \emph{Say it:} a sentence of your own with \emph{efficāx, efficācis}
  \item \emph{Recall:} say the Latin for a push, then the Latin for to bring about, then say \emph{efficāx, efficācis} again
\end{itemize}

\begin{tcolorbox}[breakable,colback=teal!4,colframe=teal!35!black,title={Before you move on}]
What does efficāx, efficācis mean? (\textquotedblleft{}Effective, efficient\textquotedblright{}.) Say it once more.
\end{tcolorbox}
\section[excitō, excitāre]{excitō, excitāre — to rouse, to wake; to stir up}

\label{lesson:LA-C372-excito}



\begin{quote}
Before the new one: say the Latin for to bring about, then the Latin for effective.
\end{quote}

\subsection*{You'll want to know: excitō, excitāre}
\textbf{excitō, excitāre} — \textquotedblleft{}to rouse, to wake; to stir up\textquotedblright{}.

\begin{quote}
Clāmor mē ē somnō excitāvit.
\end{quote}

\begin{quote}
\emph{The shouting woke me from sleep.}
\end{quote}

\emph{excitō, excitāre, excitāvī, excitātum}. To call out and set moving: a sleeper, a fire, a feeling. English \emph{excite}.

\begin{cousinweb}
\emph{excitō} is the frequentative of \emph{exciō} \textquotedblleft{}to call out, to rouse\textquotedblright{} (Lewis \& Short), from \emph{ex} \textquotedblleft{}out\textquotedblright{} and \emph{cieō} \textquotedblleft{}to set moving\textquotedblright{}, the verb behind \emph{citō}. English \emph{excite} comes from it.
\end{cousinweb}

\begin{grammarlens}[title={What you've built}]
One more word for telling, explaining and describing at B1.
\end{grammarlens}

\subsection*{Your turn}
\begin{itemize}
\raggedright
  \item \emph{Say it:} \emph{excitō, excitāre}
  \item \emph{Say it:} \emph{excitō, excitāre}, once more
  \item \emph{Say it:} a sentence of your own with \emph{excitō, excitāre}
  \item \emph{Recall:} say the Latin for to bring about, then the Latin for effective, then say \emph{excitō, excitāre} again
\end{itemize}

\begin{tcolorbox}[breakable,colback=teal!4,colframe=teal!35!black,title={Before you move on}]
What does excitō, excitāre mean? (\textquotedblleft{}To rouse, to wake; to stir up\textquotedblright{}.) Say it once more.
\end{tcolorbox}
\section[(dialogue)]{Words from chapters 365–372 — a review}

\label{lesson:LA-R372-review-365-372}



\begin{quote}
Say the Latin for since and of course.
\end{quote}

\subsection*{Your turn}
Say these aloud:

\begin{itemize}
\raggedright
  \item since, of course, for indeed, surely, to be sure
  \item evidently, without doubt, and yet, therefore, for that reason
  \item when, so that not, but that, that … not, as
  \item a speaker, eloquence, a discussion, to discuss, to speak out
  \item to urge, to advise against, persuasion, to excuse, an excuse
  \item authority, worth, weight, lightness, rashness
  \item the speakers' platform, a public meeting, to proclaim, to read aloud, rashly
  \item an originator, a push, to bring about, effective, to rouse
\end{itemize}

\begin{tcolorbox}[breakable,colback=teal!4,colframe=teal!35!black,title={Before you move on}]
Since? (\textbf{quandōquidem}.) Of course? (\textbf{quippe}.) For indeed? (\textbf{etenim}.) Surely? (\textbf{nempe}.) To be sure? (\textbf{scīlicet}.) Evidently? (\textbf{vidēlicet}.) Without doubt? (\textbf{nīmīrum}.) And yet? (\textbf{atquī}.) Therefore? (\textbf{proinde}.) For that reason? (\textbf{idcircō}.) When? (\textbf{cum}.) So that not? (\textbf{nē}.) But that? (\textbf{quīn}.) That … not? (\textbf{quōminus}.) As? (\textbf{tamquam}.) A speaker? (\textbf{ōrātor, ōrātōris}.) Eloquence? (\textbf{ēloquentia, ēloquentiae}.) A discussion? (\textbf{disputātiō, disputātiōnis}.) To discuss? (\textbf{disserō, disserere}.) To speak out? (\textbf{ēloquor, ēloquī}.) To urge? (\textbf{hortor, hortārī}.) To advise against? (\textbf{dissuādeō, dissuādēre}.) Persuasion? (\textbf{persuāsiō, persuāsiōnis}.) To excuse? (\textbf{excūsō, excūsāre}.) An excuse? (\textbf{excūsātiō, excūsātiōnis}.) Authority? (\textbf{auctōritās, auctōritātis}.) Worth? (\textbf{dignitās, dignitātis}.) Weight? (\textbf{gravitās, gravitātis}.) Lightness? (\textbf{levitās, levitātis}.) Rashness? (\textbf{temeritās, temeritātis}.) The speakers' platform? (\textbf{rōstra, rōstrōrum}.) A public meeting? (\textbf{cōntiō, cōntiōnis}.) To proclaim? (\textbf{prōnūntiō, prōnūntiāre}.) To read aloud? (\textbf{recitō, recitāre}.) Rashly? (\textbf{temerē}.) An originator? (\textbf{auctor, auctōris}.) A push? (\textbf{impulsus, impulsūs}.) To bring about? (\textbf{efficiō, efficere}.) Effective? (\textbf{efficāx, efficācis}.) To rouse? (\textbf{excitō, excitāre}.)
\end{tcolorbox}
